\section{Unannotated syntax ($\mathcal{T}_R$ and $\mathcal{T}_P$)}\label{sectionT2T2'}

Finally, we prove our last equivalence result between $\mathcal{T}_R$ and $\mathcal{T}_P$, relating
annotated and unannotated syntax. Concretely, the target theory $\mathcal{T}_P$ is the closest we can
get to a presentation as a Pure Type System of Martin-Löf's type theory with a cumulative hierarchy of universes.
This section is a direct adaptation of the arguments of \citet{Streicher} to our setting. Notably, we need
to adapt the arguments because the theory $\mathcal{T}_R$ does not enjoy uniqueness of types, precisely because
of cumulativity.

Just as we did for the previous equivalence result, we prove that the forgetful map $\strip$ preserves
derivations.

\begin{theorem}
    If $\Gamma \vdash J$ in $\mathcal{T}_R$, then we have a derivation $\strip(\Gamma) \vdash \strip(J)$ in $\mathcal{T}_P$.
\end{theorem}
\begin{proof}
    We just need to prove that the stripping of rules in $\mathcal{T}_R$ are admissible in $\mathcal{T}_P$, which
    is quite easy as each rule in $\mathcal{T}_R$ corresponds directly to a rule of $\mathcal{T}_P$. Then, the theorem
    is proved by an induction on the derivation which replaces rules in $\mathcal{T}_R$ by the admissible ones in $\mathcal{T}_P$.
\end{proof}

Now, following the ideas of \citet{Streicher}, we want to prove that two terms or types which
become syntactically equal after erasing the annotations with $\strip$ are convertible in $\mathcal{T}_R$. Here, 
types are not unique because of cumulativity, so we could well have $\strip(u_0) = \strip(u_1)$ without
necessarily having $u_0$ and $u_1$ convertible. This is due to the codomain annotations of $\lambda$-abstractions, which are erased by $\strip$ and are subject to cumulativity, as shown by the following counter-example.

\begin{example}\label{ex:cumul-lam}
    For $n < p$, the terms $f_0 = \lambda(\U_n,\U_n,v_0) : \Pi(\U_n,\U_n)$ and $f_1 = \lambda(\U_n,\U_p,v_0) : \Pi(\U_n,\U_p)$ have the same stripping $\lambda(\U_n,v_0)$,
    but neither their types $\Pi(\U_n,\U_n)$ and $\Pi(\U_n,\U_p)$ nor $f_0$ and $f_1$ are convertible.
    Yet, when applied to $\U_0$ (for $0 < n < p$), the applications $u_0 = \app(\U_n,\U_n,f_0,\U_0)$ and $u_1 = \app(\U_n,\U_p,f_1,\U_0)$ both have type $\U_p$ by cumulativity, have the same stripping $(\lambda(\U_n,v_0))\;\U_0$, and are convertible: $\Gamma \vdash u_0 = u_1 : \U_p$.
\end{example}

As Example \ref{ex:cumul-lam} illustrates, even if two applications $u_0$ and $u_1$ are compared at a common type, the functions $f_0$ and $f_1$ being applied may have types $\Pi(\U_n,\U_n)$ and $\Pi(\U_n,\U_p)$ that differ by cumulativity \emph{under $\Pi$}.
Although $\mathcal{T}_R$ has no primitive sub-typing judgement for $\Pi$-types, we can relate such terms by defining explicit coercions into a common upper type (a \emph{join}) using $\eta$-expansion, with cumulativity acting at the leaves.

\begin{definition}[Joins and coercions]\label{def:join}
    We inductively define the judgement $\Gamma \vdash \JoinTy(T_0, T_1, L) \leadsto (\coe_0, \coe_1)$, relating types $T_0, T_1$ to a common join $L$ via coercion operations $\coe_0, \coe_1$, by the following inference rules:
    \begin{mathparsmall}
      \ebrule{
        \hypo{\Gamma \vdash T_0 = L}
        \hypo{\Gamma \vdash T_1 = L}
        \infer2{\Gamma \vdash \JoinTy(T_0, T_1, L) \leadsto (\Id, \Id)}
      }
      \and
      \ebrule{
        \hypo{\Gamma \vdash T_0 = \U_k}
        \hypo{\Gamma \vdash T_1 = \U_m}
        \infer2{\Gamma \vdash \JoinTy(T_0, T_1, \U_{\max(k,m)}) \leadsto (\Id, \Id)}
      }
      \and
      \ebrule{
        \hypo{\Gamma \vdash T_0 = \Pi(A, B_0)}
        \hypo{\Gamma \vdash T_1 = \Pi(A, B_1)}
        \hypo{\Gamma.A \vdash \JoinTy(B_0, B_1, L') \leadsto (\coe_0, \coe_1)}
        \infer3{\Gamma \vdash \JoinTy(T_0, T_1, \Pi(A, L')) \leadsto (\pi_{A,B_0,L'}(\coe_0), \pi_{A,B_1,L'}(\coe_1))}
      }
    \end{mathparsmall}
    where the action of a coercion on a term $t$ is defined inductively by:
    \begin{align*}
        \Id(t) &:= t,\\
        \pi_{A,B,L'}(\coe)(t) &:= \lambda(A, L', \coe(\app(A\p, B\p^+, t\p, v_0))).
    \end{align*}
\end{definition}

By induction on the join derivation, joins are symmetric, transport along type conversions $\Gamma \vdash T'_i = T_i$, and are stable under substitution $\sigma : \Delta \to \Gamma$, with $(\coe_i(t))\sigma = \coe_{i\sigma}(t\sigma)$. Coercions are well-typed ($\Gamma \vdash t : T_i \implies \Gamma \vdash \coe_i(t) : L$), congruent, and compute under application by the $\beta$-rule:
\[\Gamma \vdash \app(A, L', \pi_{A,B_i,L'}(\coe_i)(t), a) = \coe_{i[a]}(\app(A, B_i, t, a)) : L'[a].\]

When $T_0$ and $T_1$ are convertible, coercions into their join are injective up to conversion (analogously to Lemma \ref{Utile}).

\begin{lemma}[Diagonal lemma]\label{lem:join-diag}
    If $\Gamma \vdash \JoinTy(T_0, T_1, L) \leadsto (\coe_0, \coe_1)$, $\Gamma \vdash t_0 : T_0$, $\Gamma \vdash t_1 : T_1$, and $\Gamma \vdash \coe_0(t_0) = \coe_1(t_1) : L$, then $\Gamma \vdash T_0 = T_1$ implies $\Gamma \vdash t_0 = t_1 : T_0$.
\end{lemma}
\begin{proof}
    By induction on $\Gamma \vdash \JoinTy(T_0, T_1, L) \leadsto (\coe_0, \coe_1)$. The conversion case is immediate.\\

    \textit{Universe case.} $\Gamma \vdash \U_k = \U_m$ gives $k = m$ by universe injectivity, so $L = \U_{\max(k,k)} = \U_k$.\\

    \textit{Product case.} $\Gamma \vdash \Pi(A, B_0) = \Pi(A, B_1)$ gives $\Gamma.A \vdash B_0 = B_1$ by $\Pi$-injectivity. Weakening $\Gamma \vdash \pi_{A,B_0,L'}(\coe_0)(t_0) = \pi_{A,B_1,L'}(\coe_1)(t_1) : \Pi(A, L')$ to $\Gamma.A$, applying both sides to $v_0$, and using the $\beta$-rule yields $\Gamma.A \vdash \coe_0(h_0) = \coe_1(h_1) : L'$ where $h_i = \app(A\p, B_i\p^+, t_i\p, v_0)$. The induction hypothesis gives $\Gamma.A \vdash h_0 = h_1 : B_0$, and congruence of $\lambda$ together with the $\eta$-rule concludes $\Gamma \vdash t_0 = t_1 : T_0$.
\end{proof}

We now state and prove the main lemma of this section by structural induction on the syntax, comparing terms at the join of their types.

\begin{lemma}\label{streichT2T2'}
    In $\mathcal{T}_R$, we have:
    \begin{enumerate}
        \item If $\Gamma \vdash u_0 : T_0$ and $\Gamma \vdash u_1 : T_1$ with $\strip(u_0) = \strip(u_1)$, then there exists a join $\Gamma \vdash \JoinTy(T_0, T_1, L) \leadsto (\coe_0, \coe_1)$ such that $\Gamma \vdash \coe_0(u_0) = \coe_1(u_1) : L$.
        \item If $\Gamma \vdash A_0 \; \type$ and $\Gamma \vdash A_1 \; \type$ with $\strip(A_0) = \strip(A_1)$, then $\Gamma \vdash A_0 = A_1$.
    \end{enumerate}
\end{lemma}

\begin{proof}
    By mutual structural induction on $u_0$ and $A_0$. Statement (2) follows directly from statement (1), since for $\Gamma \vdash A_i : \U_{l_i}$, any join of $\U_{l_0}$ and $\U_{l_1}$ gives $\Gamma \vdash A_0 = A_1 : \U_{\max(l_0, l_1)}$ by no-confusion and universe injectivity.
    For statement (1), inversion on $\Gamma \vdash u_i : T_i$ gives the principal type $P_i$ of $u_i$, which is related to $T_i$ either by conversion or by universe cumulativity $\U_k \leq \U_m$. It therefore suffices to construct a join for the principal types $P_0, P_1$.

    Since $\strip(u_0) = \strip(u_1)$, $u_0$ and $u_1$ have the same head constructor. Variables, universes, and products are straightforward. For $\lambda$-abstractions $u_i = \lambda(A_i, B_i, b_i)$, statement (2) gives $\Gamma \vdash A_0 = A_1$. After transporting $b_1 : B_1$ to context $\Gamma.A_0$, the induction hypothesis on $b_0, b_1$ yields a join $\Gamma.A_0 \vdash \JoinTy(B_0, B_1, L') \leadsto (\coe_0, \coe_1)$ with $\Gamma.A_0 \vdash \coe_0(b_0) = \coe_1(b_1) : L'$, which lifts to a product join on $\Pi(A_0, L')$ via the $\beta$-rule under $\lambda$. We detail the remaining case of function application.\\

    \textit{Function application.} If $u_0 = \app(C_0, D_0, c_0, a_0)$, then $u_1 = \app(C_1, D_1, c_1, a_1)$ with $\strip(c_0) = \strip(c_1)$ and $\strip(a_0) = \strip(a_1)$. By inversion, $\Gamma \vdash c_i : \Pi(C_i, D_i)$ and $\Gamma \vdash a_i : C_i$, with principal type $P_i = D_i[a_i]$.
    By the induction hypothesis, we have a join $\Gamma \vdash \JoinTy(\Pi(C_0, D_0), \Pi(C_1, D_1), L) \leadsto (\coe_0, \coe_1)$ with $\Gamma \vdash \coe_0(c_0) = \coe_1(c_1) : L$, and a join for $a_0, a_1$ on $C_0, C_1$. By no-confusion, the join for $c_0, c_1$ is either a conversion join or a product join:
    \begin{itemize}
        \item \underline{Conversion join:} We have $\Gamma \vdash \Pi(C_0, D_0) = \Pi(C_1, D_1)$ and $\Gamma \vdash c_0 = c_1 : \Pi(C_0, D_0)$. By $\Pi$-injectivity, $\Gamma \vdash C_0 = C_1$ and $\Gamma.C_0 \vdash D_0 = D_1$. Lemma \ref{lem:join-diag} on the join for $a_0, a_1$ gives $\Gamma \vdash a_0 = a_1 : C_0$, and congruence of $\app$ concludes.
        \item \underline{Product join:} We have $L = \Pi(A, L')$ with $\Gamma \vdash \Pi(C_i, D_i) = \Pi(A, B_i)$ and $\Gamma.A \vdash \JoinTy(B_0, B_1, L') \leadsto (\coe_{B_0}, \coe_{B_1})$. By $\Pi$-injectivity, $\Gamma \vdash C_i = A$ and $\Gamma.C_i \vdash D_i = B_i$. In particular, $\Gamma \vdash C_0 = C_1$, so Lemma \ref{lem:join-diag} again gives $\Gamma \vdash a_0 = a_1 : A$. Applying $\Gamma \vdash \coe_0(c_0) = \coe_1(c_1) : \Pi(A, L')$ to $a_0$, using the $\beta$-computation rule for coercions on both sides, and congruence of $\app$ yields:
        \[\Gamma \vdash \coe_{B_0[a_0]}(u_0) = \coe_{B_1[a_0]}(u_1) : L'[a_0],\]
        which is a join for $B_0[a_0]$ and $B_1[a_0]$, and hence for $D_0[a_0]$ and $D_1[a_1]$.
    \end{itemize}
\end{proof}

Combining Lemma \ref{streichT2T2'} with the diagonal lemma (Lemma \ref{lem:join-diag}), we deduce the following direct corollary, analogous to Corollary \ref{uniqueness} of the previous section.

\begin{corollary}\label{cor:strip-uniq}
    If $\Gamma \vdash u_0 : A_0$ and $\Gamma \vdash u_1 : A_1$ in $\mathcal{T}_R$ with $\strip(u_0) = \strip(u_1)$ and $\strip(A_0) = \strip(A_1)$, then $\Gamma \vdash A_0 = A_1$ and $\Gamma \vdash u_0 = u_1 : A_0$.
\end{corollary}
\begin{proof}
    If $\strip(A_0) = \strip(A_1)$, Lemma \ref{streichT2T2'}.(2) gives $\Gamma \vdash A_0 = A_1$. Then Lemma \ref{streichT2T2'}.(1) and Lemma \ref{lem:join-diag} yield $\Gamma \vdash u_0 = u_1 : A_0$.
\end{proof}

Finally, we can show the following theorem, which exhibits the section of the forgetful map $\strip$, showing the equivalence
of $\mathcal{T}_R$ and $\mathcal{T}_P$.

\begin{theorem}\label{equT2T2'}
    If $\Gamma \vdash J$ in $\mathcal{T}_P$, then we have a unique up to conversion judgment 
    $\Gamma' \vdash J'$ in $\mathcal{T}_R$ such that $\strip(\Gamma') = \Gamma$ and $\strip(J') = J$.
\end{theorem}
\begin{proof}
    By induction on the derivation $\Gamma \vdash J$, using Corollary \ref{cor:strip-uniq}, we construct a derivation in $\mathcal{T}_R$.
    The main point is that the typing annotations can be inferred from the derivation, and that Corollary \ref{cor:strip-uniq} gives
    us uniqueness up to conversion for rules with shared terms or types in their premises. 
\end{proof}

In $\mathcal{T}_P$, we could also remove the judgement $\Gamma \vdash A\; \type$ to manipulate only universes,
in a completely Russellian style. In such a system, there is no top-level of types and a type is simply a term
of a certain universe.
We can also consider named variables instead of de Bruijn indices, to get a theory presented exactly like a pure type
system, with the additional rule for cumulativity.

If we wanted to go even further, we could also remove the annotation of the variable from the lambda abstraction, again
following \citet{Streicher}. However to do so, we crucially need to rely on the full normalisation result,
as removing this annotation can lead to multiple non-convertible derivations for the same judgement, as shown by the following example.
    
\begin{example}
        We have two derivations : $t :\U_n \vdash \lambda \,x.x : \Pi(\U_n,\U_n)$ and $t :\U_n \vdash \lambda \,x.x : \Pi(\U_n,\U_p)$, which are not convertible when $n<p$.
        Then we can derive the same judgement from those two premises: \[t :\U_n \vdash (\lambda \,y.t) \; (\lambda \,x.x) : \U_n\]
\end{example}
